\section{Organizational structure}
	Contents of \textbf{~/porous-media/}
	\begin{enumerate}
		\item materials/: contains articles, books, videos.
		\item network-model/: the code and its documentation.
		\item work-record/: old results and calculations.
	\end{enumerate}

	Contents of \textbf{~/porous-media/network-model/}:
	\begin{enumerate}
		\item makefilegen/: this is the makefile generator program.
		\item manual/: documentation.
		\item run/: temporary files and results.
		\item src/: the source code.
		\item Makefile: scripts to ease compilation and running of the program.
		\item README.md: information for git hub.
	\end{enumerate}


\section{makefilegen program}
	Makefilegen is a program. This program is an exe file, compiled from a C++ source code. It reads makefilegen-config.txt which cointains information about the sourcecode and generates a Makefile. This makefile eases the process of compilation.

	This must be the file and folder structure \textbf{network-model/} for makefilegen to function properly:
	\begin{enumerate}
		\item makefilegen/
			\begin{enumerate}
				\item src/
					\begin{enumerate}
						\item decl.h: basic declaration file containing basic file and folder names
						\item main.cpp: contains int main(), uses the Class Menu, and runs the most basic commands from Class Menu, all of which writes some
						\item menu.cpp
						\item menu.h
						\item print.cpp
						\item print.h
					\end{enumerate}
				\item build/
		\item manual/: documentation.
		\item run/: temporary files and results.
		\item src/: the source code.
		\item Makefile: scripts to ease compilation and running of the program.
		\item README.md: information for git hub.
	\end{enumerate}
	\end{enumerate}
